% Generated by publication/build_sections.py; do not edit by hand.
\paragraph{Informal verdict.}
\textbf{A --- Complete solution}, with \textbf{high confidence}.

\paragraph{Lean coverage.}
Lean proves the full arbitrary-index compact-topological-group translation-separation theorem without metrizability or countability assumptions.

\subsection{Audited informal answer}
Q781 --- Complete solution
\begin{enumerate}
\item Faithful restatement and interpretation
\end{enumerate}
Let \(G\) be a compact topological group, and let \((A_i)_{i\in I}\) be a family of subsets of \(G\). The family is called separable by translations if there are elements \(t_i\in G\) such that the left translates \(t_iA_i\), \(i\in I\), are pairwise disjoint.
Given a family \((U_i)_{i\in I}\) of open subsets of \(G\), assume that every finite subfamily is separable by translations. Prove that the whole family is separable by translations.
The PDF first says merely ``let \(G\) be a group'' and then speaks of compactness and open subsets. The intended meaning must therefore be ``compact topological group.'' The word ``separable'' is the specially defined translation-packing property, not separability in the sense of possessing a countable dense subset. The parenthetical remark permits assuming metrizability and countability if needed, but neither is necessary. There are no further subparts or footnotes. Luc Pommeret+1

\begin{enumerate}
\item The theorem
\end{enumerate}
Theorem
Let \(G\) be a compact topological group, let \(I\) be an arbitrary index set, and let \(U_i\subseteq G\) be open for every \(i\in I\). Suppose that for every finite \(F\subseteq I\), there exist elements \(t_i^{(F)}\in G\), \(i\in F\), such that

\[
t_i^{(F)}U_i\cap t_j^{(F)}U_j=\varnothing
\qquad(i,j\in F,\ i\neq j).
\]

Then there exist elements \(t_i\in G\), \(i\in I\), such that

\[
t_iU_i\cap t_jU_j=\varnothing
\qquad(i,j\in I,\ i\neq j).
\]

Thus Q781 has an affirmative answer in its full stated generality: neither metrizability of \(G\) nor countability of \(I\) is needed.

\begin{enumerate}
\item The disjointness condition is closed
\end{enumerate}
The crucial point is that, because the \(U_i\) are open, collision of two translated sets is an open condition on the two translation parameters.
Fix distinct \(i,j\in I\). Define

\[
R_{ij}
=
\bigl\{(s,t)\in G\times G:\ sU_i\cap tU_j=\varnothing\bigr\}.
\]

We claim that \(R_{ij}\) is closed in \(G\times G\).
Indeed, for \(s,t\in G\),

\[
\begin{aligned}
sU_i\cap tU_j\neq\varnothing
&\iff
\exists\,u\in U_i,\ v\in U_j\text{ such that }su=tv\\
&\iff
\exists\,u\in U_i,\ v\in U_j\text{ such that }s^{-1}t=uv^{-1}\\
&\iff
s^{-1}t\in U_iU_j^{-1}.
\end{aligned}
\]

The subset \(U_iU_j^{-1}\) is open. To see this, inversion is a homeomorphism, so \(U_j^{-1}\) is open, and

\[
U_iU_j^{-1}
=
\bigcup_{u\in U_i}uU_j^{-1}
\]

is a union of open left translates.
The map

\[
q:G\times G\longrightarrow G,
\qquad
q(s,t)=s^{-1}t,
\]

is continuous. Consequently,

\[
(G\times G)\setminus R_{ij}
=
q^{-1}(U_iU_j^{-1})
\]

is open. Hence \(R_{ij}\) is closed.
Notice the useful direction of semicontinuity here: a limit of pairs of translations whose corresponding open sets are disjoint still gives disjoint open sets. In contrast, a limit of intersecting translated open sets need not remain intersecting.

\begin{enumerate}
\item Compactness of the full parameter space
\end{enumerate}
Consider the product space

\[
X=G^I=\prod_{i\in I}G
\]

with the product topology. A point \(x\in X\) will be written

\[
x=(x_i)_{i\in I}.
\]

By Tychonoff's theorem, \(X\) is compact.
For every unordered pair \(\{i,j\}\subseteq I\), \(i\neq j\), define

\[
K_{ij}
=
\bigl\{x=(x_k)_{k\in I}\in X:
x_iU_i\cap x_jU_j=\varnothing\bigr\}.
\]

Let

\[
p_{ij}:X\longrightarrow G\times G,
\qquad
p_{ij}(x)=(x_i,x_j)
\]

be the coordinate projection. Then

\[
K_{ij}=p_{ij}^{-1}(R_{ij}).
\]

Since \(R_{ij}\) is closed and \(p_{ij}\) is continuous, \(K_{ij}\) is closed in \(X\).
We now prove that the family

\[
\{K_{ij}:i,j\in I,\ i\neq j\}
\]

has the finite-intersection property.
Take finitely many pair constraints,

\[
K_{i_1j_1},\ldots,K_{i_mj_m}.
\]

Let

\[
F=\{i_1,j_1,\ldots,i_m,j_m\}\subseteq I.
\]

This is a finite set of indices. By hypothesis, the finite family \((U_i)_{i\in F}\) can be separated: there exist \(s_i\in G\), \(i\in F\), such that

\[
s_iU_i\cap s_jU_j=\varnothing
\qquad(i,j\in F,\ i\neq j).
\]

Extend this finite assignment to a point \(x\in G^I\) by setting

\[
x_i=
\begin{cases}
s_i,&i\in F,\\
e,&i\notin F,
\end{cases}
\]

where \(e\) is the identity of \(G\). Then

\[
x\in K_{i_1j_1}\cap\cdots\cap K_{i_mj_m}.
\]

Thus every finite intersection of the \(K_{ij}\) is nonempty.
Because \(X\) is compact and the \(K_{ij}\) are closed with the finite-intersection property,

\[
\bigcap_{\substack{i,j\in I\\\i\neq j}}K_{ij}\neq\varnothing.
\]

Choose

\[
(t_i)_{i\in I}
\in
\bigcap_{i\neq j}K_{ij}.
\]

Then, by the definition of \(K_{ij}\),

\[
t_iU_i\cap t_jU_j=\varnothing
\qquad(i\neq j).
\]

Therefore the entire family \((U_i)_{i\in I}\) is separable by translations. \(\square\)

\begin{enumerate}
\item The countable metrizable proof
\end{enumerate}
The optional assumptions in the PDF allow a purely sequential formulation.
Suppose \(I=\mathbb N\) and \(G\) is compact metrizable. For each \(n\geq 1\), choose translations

\[
t_1^{(n)},\ldots,t_n^{(n)}
\]

separating \(U_1,\ldots,U_n\), and extend this to a point \(x^{(n)}\in G^{\mathbb N}\) by setting

\[
x_k^{(n)}=
\begin{cases}
t_k^{(n)},&k\leq n,\\
e,&k>n.
\end{cases}
\]

If \(d\) is a compatible metric on \(G\), the product topology on \(G^{\mathbb N}\) is induced by, for example,

\[
D(x,y)
=
\sum_{k=1}^{\infty}2^{-k}\min\{1,d(x_k,y_k)\}.
\]

The countable product \(G^{\mathbb N}\) is compact metrizable. Equivalently, one may obtain a convergent subsequence by successively extracting subsequences on which the first coordinate converges, then the second, and so on, and taking the diagonal subsequence.
Thus, after passing to a subsequence,

\[
x^{(n_r)}\longrightarrow x=(t_k)_{k\geq1}
\quad\text{in }G^{\mathbb N}.
\]

Fix \(i\neq j\). Since \(n_r\to\infty\), for every sufficiently large \(r\),

\[
x_i^{(n_r)}U_i
\cap
x_j^{(n_r)}U_j
=
\varnothing.
\]

The relation \(R_{ij}\) proved closed above, so passage to the limit gives

\[
t_iU_i\cap t_jU_j=\varnothing.
\]

Since \(i,j\) were arbitrary, the limit tuple separates the entire family.
Why this is not the general proof
For a general compact group, compactness is not a sequential notion. A compact space need not be sequentially compact, and an uncountable product \(G^I\) generally need not be metrizable or first countable.
For example, the compact Boolean group

\[
H=\{0,1\}^{\mathcal P(\mathbb N)}
\]

is not sequentially compact. Define \(x_n\in H\) by

\[
x_n(A)=
\begin{cases}
1,&n\in A,\\
0,&n\notin A.
\end{cases}
\]

Given any subsequence \((x_{n_k})\), take

\[
A=\{n_{2k}:k\geq1\}.
\]

Then the \(A\)-coordinate of \(x_{n_k}\) alternates between \(0\) and \(1\), so the subsequence does not converge. Thus a proof that simply says ``extract a convergent subsequence'' would be invalid without the metrizability assumptions. The finite-intersection argument above is genuinely non-sequential and works in the full generality of Q781.

\begin{enumerate}
\item Audit of the topological and set-theoretic assumptions
\end{enumerate}
Separation assumptions
No Hausdorff, regularity, normality, or first-countability axiom is used in the argument itself. The proof only uses:


continuity of multiplication and inversion;


openness of the \(U_i\);


compactness of \(G^I\).


Under the usual convention, a ``compact group'' is a compact Hausdorff topological group, but Hausdorffness is not needed for the closed-constraint argument.
Compactness
The general proof uses the arbitrary-product form of Tychonoff's theorem:

The product of an arbitrary family of compact spaces is compact in the product topology.

Compactness is then used through its equivalent finite-intersection formulation: a family of closed subsets of a compact space whose finite intersections are all nonempty has nonempty total intersection.
Countability
No countability assumption on \(I\) is used. No second-countability or separability assumption on \(G\) is used.
Countability becomes relevant only if one wishes to replace the product-compactness proof by the sequential diagonal argument.
Choice
In ordinary ZFC, Tychonoff's theorem is available. Foundationally, for compact Hausdorff spaces, the ultrafilter lemma---equivalently the Boolean prime ideal theorem---suffices for the required product compactness; full choice is more than is needed for that step.
No choice is needed to extend a finite tuple to all of \(I\): the identity \(e\) supplies a canonical value on every unused coordinate.
In the countable sequential proof, literally choosing one finite separator for every \(n\) uses a countable selection principle, which is again automatic in ZFC.
The precise role of openness
Openness is used exactly to prove that

\[
\{(s,t):sU_i\cap tU_j=\varnothing\}
\]

is closed. In fact, the same proof works under the weaker pairwise assumption

\[
U_iU_j^{-1}\ \text{is open for every }i\neq j.
\]

Thus openness of each individual \(U_i\) is sufficient but is not the logically weakest possible condition.

\begin{enumerate}
\item Boundary cases
\end{enumerate}
If \(I=\varnothing\) or \(|I|=1\), there are no pairwise constraints, so the conclusion is immediate.
If some \(U_i=\varnothing\), that index causes no constraint because every translate of \(U_i\) is empty. Such indices can be ignored and then assigned any translation, for example \(e\).
If \(G\) is a finite discrete group of order \(N\), the hypothesis already implies that there are at most \(N\) nonempty \(U_i\). Indeed, \(N+1\) pairwise disjoint nonempty subsets cannot exist inside \(G\). The set of nonempty indices is therefore finite, and applying the hypothesis to that finite subfamily immediately gives the conclusion. This agrees with the general compactness proof.

\begin{enumerate}
\item Why compactness cannot be omitted
\end{enumerate}
The assertion fails for a noncompact group even when:


the group is discrete, metrizable, locally compact, and countable;


\(I\) is countable;


every \(U_i\) is open.


Take

\[
G=(\mathbb Z,+)
\]

with the discrete topology. Let the index set be

\[
I=\{L,R\}\cup\mathbb N,
\]

and define

\[
U_L=\{m\in\mathbb Z:m\leq0\},\qquad
U_R=\{m\in\mathbb Z:m\geq0\},
\]

and

\[
U_n=\{0\}
\qquad(n\in\mathbb N).
\]

Every finite subfamily can be separated. Indeed, suppose it contains \(k\) of the singleton sets, indexed in some order by \(n_1,\ldots,n_k\). Use the translations

\[
t_L=0,\qquad t_R=k+1,\qquad t_{n_r}=r
\quad(1\leq r\leq k).
\]

Then

\[
t_L+U_L=(-\infty,0]\cap\mathbb Z,
\]


\[
t_R+U_R=[k+1,\infty)\cap\mathbb Z,
\]

and the singleton translates are

\[
\{1\},\ldots,\{k\}.
\]

These sets are pairwise disjoint. The same assignment works if one or both of \(L,R\) are absent from the finite subfamily.
However, the entire family cannot be separated. Suppose translations \(a,b,c_n\in\mathbb Z\) existed. Disjointness of the two rays would require

\[
a<b.
\]

For the singleton \(\{c_n\}\) to avoid both rays, one must have

\[
a<c_n<b.
\]

Moreover, the \(c_n\) must be pairwise distinct. But the finite interval

\[
\{a+1,\ldots,b-1\}
\]

contains only finitely many integers. It cannot contain all the distinct \(c_n\). This contradiction proves that compactness is essential for a theorem of this generality.

\begin{enumerate}
\item Why arbitrary non-open subsets cannot be allowed
\end{enumerate}
The result also fails for general subsets of a compact group, even if all the subsets are closed.
Let

\[
G=\mathbb T=\mathbb R/\mathbb Z
\]

and choose an index set \(I\) with cardinality strictly greater than \(|\mathbb T|\), for instance

\[
I=\mathcal P(\mathbb T).
\]

For every \(i\in I\), put

\[
A_i=\{0\}.
\]

Every finite subfamily is separable: simply choose finitely many distinct elements \(t_i\in\mathbb T\).
A separation of the entire family would require the singleton translates

\[
t_iA_i=\{t_i\},
\qquad i\in I,
\]

to be pairwise disjoint. Thus \(i\mapsto t_i\) would be an injection from \(I\) into \(\mathbb T\), which is impossible because

\[
|I|=|\mathcal P(\mathbb T)|>|\mathbb T|.
\]

Here \(G\) is compact Hausdorff and every \(A_i\) is closed, but the \(A_i\) are not open. Hence the openness/closed-collision condition cannot simply be discarded.

Final conclusion
The answer to Q781 is yes. For each pair \(i\neq j\), disjointness of \(t_iU_i\) and \(t_jU_j\) is a closed condition on the parameter tuple \((t_k)_{k\in I}\in G^I\). The finite-separability hypothesis says that these closed conditions have the finite-intersection property. Tychonoff compactness of \(G^I\) therefore supplies one tuple satisfying all pairwise conditions simultaneously.
